% =====================================================================
%  CONJURA CONJECTURE STATEMENT -- TEMPLATE
%
%  Copy this whole folder (conjura-conjecture.cls, conjura-solution.cls,
%  statement.tex, solution.tex) to start a new conjecture.  Fill in
%  every placeholder below and delete this comment block.
%
%  This document is the STATEMENT only.  A proof or attempted proof,
%  if and when one exists, is a separate document -- see solution.tex,
%  not this file.
% =====================================================================
\documentclass{conjura-conjecture}

\runninghead{SHORT TITLE IN CAPS}

\begin{document}

\cjkicker{CONJURA \textperiodcentered{} OPEN PROBLEM}
\cjtitle{Conjecture Title Here}
\cjsubtitle{One-line subtitle, if useful}
\cjstatus{Statement: AI-written, not yet formalized.}
\cjcategory{Fill in the site category here, e.g.\ \emph{Proof Systems
  \& Succinct Arguments}.}

\vspace{0.4em}

\begin{informalconjecture}
One paragraph, in plain English, informally restating what the
conjecture claims.
\end{informalconjecture}

\section{The setting}\label{sec:setting}

Background and notation, where the problem originates, related
work, and progress to date.

\section{The conjectures}\label{sec:conjectures}

Formal, self-contained definitions, followed by the conjecture(s)
themselves.

\begin{definition}[name]\label{def:example}
State any definitions the conjecture needs here.
\end{definition}

\begin{conjecture}[name]\label{conj:main}
State the conjecture itself here.
\end{conjecture}

\section{Bibliography}

\begin{conjurabibliography}{99}

\bibitem[Key]{Key}
Author.
\emph{Title}.
Venue, year.

\end{conjurabibliography}

\end{document}
